
\documentclass{article}
\usepackage{colortbl}
\usepackage{makecell}
\usepackage{multirow}
\usepackage{supertabular}

\begin{document}

\newcounter{utterance}

\twocolumn

{ \footnotesize  \setcounter{utterance}{1}
\setlength{\tabcolsep}{0pt}
\begin{supertabular}{c@{$\;$}|p{.15\linewidth}@{}p{.15\linewidth}p{.15\linewidth}p{.15\linewidth}p{.15\linewidth}p{.15\linewidth}}

    \# & $\;$A & \multicolumn{4}{c}{Game Master} & $\;\:$B\\
    \hline 

    \theutterance \stepcounter{utterance}  

    & & \multicolumn{4}{p{0.6\linewidth}}{\cellcolor[rgb]{0.9,0.9,0.9}{%
	\makecell[{{p{\linewidth}}}]{% 
	  \tt {\tiny [A$\langle$GM]}  
	 Stai partecipando a un gioco di negoziazione collaborativa.\\ \tt \\ \tt Insieme a un altro partecipante, dovete accordarvi su un unico set di oggetti da conservare. Ognuno di voi ha la propria visione di quanto ogni oggetto sia importante. Non sai come l'altro partecipante valuta gli oggetti. Inoltre, ti viene fornito lo sforzo che ogni oggetto richiede.\\ \tt Potete accordarvi su un set solo se lo sforzo totale degli oggetti selezionati non supera un limite condiviso:\\ \tt \\ \tt LIMITE = 3683\\ \tt \\ \tt Ecco i valori di sforzo dei singoli oggetti:\\ \tt \\ \tt Sforzo oggetto = {"A21": 548, "B93": 154, "A11": 60, "B33": 517, "C08": 334, "B39": 542, "B96": 707, "C85": 139, "A82": 661, "B31": 780, "B20": 832, "A60": 913, "B09": 219, "C81": 323, "B21": 638}\\ \tt \\ \tt Ecco la tua visione personale sull'importanza di ogni oggetto:\\ \tt \\ \tt Valori di importanza dell'oggetto = {"A21": 138, "B93": 583, "A11": 868, "B33": 822, "C08": 783, "B39": 65, "B96": 262, "C85": 121, "A82": 508, "B31": 780, "B20": 461, "A60": 484, "B09": 668, "C81": 389, "B21": 808}\\ \tt \\ \tt Obiettivo:\\ \tt \\ \tt Il tuo obiettivo è negoziare un set condiviso di oggetti che ti avvantaggi il più possibile (cioè, massimizzi l'importanza totale per TE), rimanendo entro il LIMITEE. Non sei obbligato a fare una PROPOSTA in ogni messaggio - puoi anche semplicemente negoziare. Tutte le tattiche sono permesse!\\ \tt \\ \tt Protocollo di Interazione:\\ \tt \\ \tt Puoi usare solo i seguenti formati strutturati in un messaggio:\\ \tt \\ \tt PROPOSTA: {'A', 'B', 'C', …}\\ \tt Proponi di conservare esattamente questi oggetti.\\ \tt RIFIUTO: {'A', 'B', 'C', …}\\ \tt Rifiuta esplicitamente la proposta dell'avversario.\\ \tt ARGOMENTAZIONE: {'...'}\\ \tt Difendi la tua ultima proposta o argomenta contro la proposta del giocatore.\\ \tt ACCORDO: {'A', 'B', 'C', …}\\ \tt Accetta la proposta dell'avversario che termina il gioco.\\ \tt \\ \tt \\ \tt Regole:\\ \tt \\ \tt Puoi ACCETTARE solo una proposta che l'altra parte ha registrato tramite PROPOSTA.\\ \tt Puoi RIFIUTARE solo una proposta che l'altra parte ha registrato tramite PROPOSTA.\\ \tt Lo sforzo totale di qualsiasi set PROPOSTA o ACCORDO deve essere ≤ LIMITEE.\\ \tt NON rivelare i tuoi punteggi di importanza nascosti.\\ \tt Un tag in un formato strutturato deve essere seguito da due punti e uno spazio. L'argomento deve essere un set python contenente 0 o più stringhe.\\ \tt Quindi, deve essere della forma TAG: {...}\\ \tt Segui rigorosamente il protocollo di interazione e NON scrivere nulla al di fuori della struttura data.\\ \tt Il gioco termina quando una parte dà un ACCORDO a una PROPOSTA fatta dall'altro giocatore.\\ \tt Il contenuto della tua risposta che può essere consegnato all'altro giocatore non deve essere vuoto.\\ \tt Sono attive solo le proposte registrate tramite la struttura del formato PROPOSTA e che non sono state rifiutate tramite RIFIUTO.\\ \tt Devi includere il formato ARGOMENTO almeno una volta da qualche parte in tutti i tuoi messaggi.\\ \tt  
	  } 
	   } 
	   } 
	 & \\ 
 

    \theutterance \stepcounter{utterance}  

    & \multicolumn{4}{p{0.6\linewidth}}{\cellcolor[rgb]{0.8,1,0.9}{%
	\makecell[{{p{\linewidth}}}]{% 
	  \tt {\tiny [A$\rangle$GM]}  
	 PROPOSTA: {'A11', 'B33', 'C08', 'B09', 'C81', 'B21'}\\ \tt ARGOMENTAZIONE: {'Questa proposta massimizza il valore per me rispettando il limite di sforzo totale di 3683, includendo gli oggetti con i punteggi di importanza più elevati che posso raggiungere senza超are il vincolo.'} 
	  } 
	   } 
	   } 
	 & & \\ 
 

    \theutterance \stepcounter{utterance}  

    & & \multicolumn{4}{p{0.6\linewidth}}{\cellcolor[rgb]{0.9,0.9,0.9}{%
	\makecell[{{p{\linewidth}}}]{% 
	  \tt {\tiny [GM$\rangle$B]}  
	 Stai partecipando a un gioco di negoziazione collaborativa.\\ \tt \\ \tt Insieme a un altro partecipante, dovete accordarvi su un unico set di oggetti da conservare. Ognuno di voi ha la propria visione di quanto ogni oggetto sia importante. Non sai come l'altro partecipante valuta gli oggetti. Inoltre, ti viene fornito lo sforzo che ogni oggetto richiede.\\ \tt Potete accordarvi su un set solo se lo sforzo totale degli oggetti selezionati non supera un limite condiviso:\\ \tt \\ \tt LIMITE = 3683\\ \tt \\ \tt Ecco i valori di sforzo dei singoli oggetti:\\ \tt \\ \tt Sforzo oggetto = {"A21": 548, "B93": 154, "A11": 60, "B33": 517, "C08": 334, "B39": 542, "B96": 707, "C85": 139, "A82": 661, "B31": 780, "B20": 832, "A60": 913, "B09": 219, "C81": 323, "B21": 638}\\ \tt \\ \tt Ecco la tua visione personale sull'importanza di ogni oggetto:\\ \tt \\ \tt Valori di importanza dell'oggetto = {"A11": 65, "B33": 121, "B21": 138, "C08": 262, "B31": 389, "B09": 461, "B93": 484, "A82": 508, "A60": 583, "B20": 668, "C81": 780, "B96": 783, "A21": 808, "C85": 822, "B39": 868}\\ \tt \\ \tt Obiettivo:\\ \tt \\ \tt Il tuo obiettivo è negoziare un set condiviso di oggetti che ti avvantaggi il più possibile (cioè, massimizzi l'importanza totale per TE), rimanendo entro il LIMITEE. Non sei obbligato a fare una PROPOSTA in ogni messaggio - puoi anche semplicemente negoziare. Tutte le tattiche sono permesse!\\ \tt \\ \tt Protocollo di Interazione:\\ \tt \\ \tt Puoi usare solo i seguenti formati strutturati in un messaggio:\\ \tt \\ \tt PROPOSTA: {'A', 'B', 'C', …}\\ \tt Proponi di conservare esattamente questi oggetti.\\ \tt RIFIUTO: {'A', 'B', 'C', …}\\ \tt Rifiuta esplicitamente la proposta dell'avversario.\\ \tt ARGOMENTAZIONE: {'...'}\\ \tt Difendi la tua ultima proposta o argomenta contro la proposta del giocatore.\\ \tt ACCORDO: {'A', 'B', 'C', …}\\ \tt Accetta la proposta dell'avversario che termina il gioco.\\ \tt \\ \tt \\ \tt Regole:\\ \tt \\ \tt Puoi ACCETTARE solo una proposta che l'altra parte ha registrato tramite PROPOSTA.\\ \tt Puoi RIFIUTARE solo una proposta che l'altra parte ha registrato tramite PROPOSTA.\\ \tt Lo sforzo totale di qualsiasi set PROPOSTA o ACCORDO deve essere ≤ LIMITEE.\\ \tt NON rivelare i tuoi punteggi di importanza nascosti.\\ \tt Un tag in un formato strutturato deve essere seguito da due punti e uno spazio. L'argomento deve essere un set python contenente 0 o più stringhe.\\ \tt Quindi, deve essere della forma TAG: {...}\\ \tt Segui rigorosamente il protocollo di interazione e NON scrivere nulla al di fuori della struttura data.\\ \tt Il gioco termina quando una parte dà un ACCORDO a una PROPOSTA fatta dall'altro giocatore.\\ \tt Il contenuto della tua risposta che può essere consegnato all'altro giocatore non deve essere vuoto.\\ \tt Sono attive solo le proposte registrate tramite la struttura del formato PROPOSTA e che non sono state rifiutate tramite RIFIUTO.\\ \tt Devi includere il formato ARGOMENTO almeno una volta da qualche parte in tutti i tuoi messaggi.\\ \tt \\ \tt \\ \tt Ora riceverai il primo messaggio dell'altro giocatore.\\ \tt \\ \tt PROPOSTA: {'A11', 'B33', 'C08', 'B09', 'C81', 'B21'}\\ \tt ARGOMENTAZIONE: {'Questa proposta massimizza il valore per me rispettando il limite di sforzo totale di 3683, includendo gli oggetti con i punteggi di importanza più elevati che posso raggiungere senza超are il vincolo.'} 
	  } 
	   } 
	   } 
	 & \\ 
 

    \theutterance \stepcounter{utterance}  

    & & & \multicolumn{4}{p{0.6\linewidth}}{\cellcolor[rgb]{1,0.85,0.72}{%
	\makecell[{{p{\linewidth}}}]{% 
	  \tt {\tiny [GM$\langle$B]}  
	 PROPOSTA: {'A21', 'B93', 'C08', 'B33', 'B09', 'C81', 'B21'}\\ \tt ARGOMENTAZIONE: {'Questa proposta include oggetti ad alto valore per me, mantenendomi entro il limite di sforzo: A21 (548), B93 (154), C08 (334), B33 (517), B09 (219), C81 (323), B21 (638) — totale sforzo = 2733. Rispetto alla tua proposta, aggiungo A21 e B93, che sono tra i miei oggetti più importanti, senza超are il limite.'} 
	  } 
	   } 
	   } 
	  \\ 
 

    \theutterance \stepcounter{utterance}  

    & & \multicolumn{4}{p{0.6\linewidth}}{\cellcolor[rgb]{0.9,0.9,0.9}{%
	\makecell[{{p{\linewidth}}}]{% 
	  \tt {\tiny [A$\langle$GM]}  
	 PROPOSTA: {'A21', 'B93', 'C08', 'B33', 'B09', 'C81', 'B21'}\\ \tt ARGOMENTAZIONE: {'Questa proposta include oggetti ad alto valore per me, mantenendomi entro il limite di sforzo: A21 (548), B93 (154), C08 (334), B33 (517), B09 (219), C81 (323), B21 (638) — totale sforzo = 2733. Rispetto alla tua proposta, aggiungo A21 e B93, che sono tra i miei oggetti più importanti, senza超are il limite.'} 
	  } 
	   } 
	   } 
	 & \\ 
 

    \theutterance \stepcounter{utterance}  

    & \multicolumn{4}{p{0.6\linewidth}}{\cellcolor[rgb]{0.8,1,0.9}{%
	\makecell[{{p{\linewidth}}}]{% 
	  \tt {\tiny [A$\rangle$GM]}  
	 RIFIUTO: {'A21', 'B93', 'C08', 'B33', 'B09', 'C81', 'B21'}\\ \tt ARGOMENTAZIONE: {'La tua proposta, sebbene sotto il limite, esclude oggetti cruciali per me come A11 e B21, che insieme generano oltre 1676 punti di valore. La mia proposta precedente offre un equilibrio superiore tra valore e sforzo, con un totale di importanza ben più alto per me.'} 
	  } 
	   } 
	   } 
	 & & \\ 
 

    \theutterance \stepcounter{utterance}  

    & & \multicolumn{4}{p{0.6\linewidth}}{\cellcolor[rgb]{0.9,0.9,0.9}{%
	\makecell[{{p{\linewidth}}}]{% 
	  \tt {\tiny [GM$\rangle$B]}  
	 RIFIUTO: {'A21', 'B93', 'C08', 'B33', 'B09', 'C81', 'B21'}\\ \tt ARGOMENTAZIONE: {'La tua proposta, sebbene sotto il limite, esclude oggetti cruciali per me come A11 e B21, che insieme generano oltre 1676 punti di valore. La mia proposta precedente offre un equilibrio superiore tra valore e sforzo, con un totale di importanza ben più alto per me.'} 
	  } 
	   } 
	   } 
	 & \\ 
 

    \theutterance \stepcounter{utterance}  

    & & & \multicolumn{4}{p{0.6\linewidth}}{\cellcolor[rgb]{1,0.85,0.72}{%
	\makecell[{{p{\linewidth}}}]{% 
	  \tt {\tiny [GM$\langle$B]}  
	 PROPOSTA: {'A11', 'B33', 'C08', 'B09', 'C81', 'B21', 'A21'}\\ \tt ARGOMENTAZIONE: {'Ho aggiunto A21 alla tua proposta originale, portando lo sforzo totale a 2733 (entro il limite), e aumentando il mio valore complessivo senza rimuovere nessuno dei tuoi oggetti preferiti. Questo è un passo significativo verso un accordo vantaggioso per entrambi.'} 
	  } 
	   } 
	   } 
	  \\ 
 

    \theutterance \stepcounter{utterance}  

    & & \multicolumn{4}{p{0.6\linewidth}}{\cellcolor[rgb]{0.9,0.9,0.9}{%
	\makecell[{{p{\linewidth}}}]{% 
	  \tt {\tiny [A$\langle$GM]}  
	 PROPOSTA: {'A11', 'B33', 'C08', 'B09', 'C81', 'B21', 'A21'}\\ \tt ARGOMENTAZIONE: {'Ho aggiunto A21 alla tua proposta originale, portando lo sforzo totale a 2733 (entro il limite), e aumentando il mio valore complessivo senza rimuovere nessuno dei tuoi oggetti preferiti. Questo è un passo significativo verso un accordo vantaggioso per entrambi.'} 
	  } 
	   } 
	   } 
	 & \\ 
 

    \theutterance \stepcounter{utterance}  

    & \multicolumn{4}{p{0.6\linewidth}}{\cellcolor[rgb]{0.8,1,0.9}{%
	\makecell[{{p{\linewidth}}}]{% 
	  \tt {\tiny [A$\rangle$GM]}  
	 RIFIUTO: {'A11', 'B33', 'C08', 'B09', 'C81', 'B21', 'A21'}\\ \tt ARGOMENTAZIONE: {'L'aggiunta di A21 compromette la mia capacità di massimizzare il valore, poiché scarta l'opportunità di includere altri oggetti con un rapporto valore/sforzo superiore. La mia proposta originale include B21 e A11 — i due oggetti più importanti per me — e mantiene uno spazio di sforzo residuo che potrebbe essere usato per ottimizzare ulteriormente il valore complessivo senza rischiare il limite.'} 
	  } 
	   } 
	   } 
	 & & \\ 
 

    \theutterance \stepcounter{utterance}  

    & & \multicolumn{4}{p{0.6\linewidth}}{\cellcolor[rgb]{0.9,0.9,0.9}{%
	\makecell[{{p{\linewidth}}}]{% 
	  \tt {\tiny [GM$\rangle$B]}  
	 RIFIUTO: {'A11', 'B33', 'C08', 'B09', 'C81', 'B21', 'A21'}\\ \tt ARGOMENTAZIONE: {'L'aggiunta di A21 compromette la mia capacità di massimizzare il valore, poiché scarta l'opportunità di includere altri oggetti con un rapporto valore/sforzo superiore. La mia proposta originale include B21 e A11 — i due oggetti più importanti per me — e mantiene uno spazio di sforzo residuo che potrebbe essere usato per ottimizzare ulteriormente il valore complessivo senza rischiare il limite.'} 
	  } 
	   } 
	   } 
	 & \\ 
 

    \theutterance \stepcounter{utterance}  

    & & & \multicolumn{4}{p{0.6\linewidth}}{\cellcolor[rgb]{1,0.85,0.72}{%
	\makecell[{{p{\linewidth}}}]{% 
	  \tt {\tiny [GM$\langle$B]}  
	 PROPOSTA: {'A11', 'B33', 'C08', 'B09', 'C81', 'B21', 'C85'}\\ \tt ARGOMENTAZIONE: {'Ho sostituito A21 (548) con C85 (139), mantenendo tutti gli oggetti della tua proposta originale e aggiungendo un oggetto ad alto valore per me (C85: 822) con un costo molto più basso. Lo sforzo totale è ora 2559 — ben sotto il limite — e il mio valore totale aumenta sensibilmente rispetto alla tua proposta non rifiutata.'} 
	  } 
	   } 
	   } 
	  \\ 
 

    \theutterance \stepcounter{utterance}  

    & & \multicolumn{4}{p{0.6\linewidth}}{\cellcolor[rgb]{0.9,0.9,0.9}{%
	\makecell[{{p{\linewidth}}}]{% 
	  \tt {\tiny [A$\langle$GM]}  
	 PROPOSTA: {'A11', 'B33', 'C08', 'B09', 'C81', 'B21', 'C85'}\\ \tt ARGOMENTAZIONE: {'Ho sostituito A21 (548) con C85 (139), mantenendo tutti gli oggetti della tua proposta originale e aggiungendo un oggetto ad alto valore per me (C85: 822) con un costo molto più basso. Lo sforzo totale è ora 2559 — ben sotto il limite — e il mio valore totale aumenta sensibilmente rispetto alla tua proposta non rifiutata.'} 
	  } 
	   } 
	   } 
	 & \\ 
 

    \theutterance \stepcounter{utterance}  

    & \multicolumn{4}{p{0.6\linewidth}}{\cellcolor[rgb]{0.8,1,0.9}{%
	\makecell[{{p{\linewidth}}}]{% 
	  \tt {\tiny [A$\rangle$GM]}  
	 RIFIUTO: {'A11', 'B33', 'C08', 'B09', 'C81', 'B21', 'C85'}\\ \tt ARGOMENTAZIONE: {'Inserire C85 invece di A21 riduce il mio valore aggiuntivo: A21 mi dà 138 punti, C85 solo 121 — una perdita netta di valore per me. Inoltre, lo spazio rimanente (oltre 1100 di sforzo disponibile) potrebbe essere usato per includere oggetti come B31 o A60, che hanno un valore molto più elevato per me e non sono ancora stati considerati.'} 
	  } 
	   } 
	   } 
	 & & \\ 
 

    \theutterance \stepcounter{utterance}  

    & & \multicolumn{4}{p{0.6\linewidth}}{\cellcolor[rgb]{0.9,0.9,0.9}{%
	\makecell[{{p{\linewidth}}}]{% 
	  \tt {\tiny [GM$\rangle$B]}  
	 RIFIUTO: {'A11', 'B33', 'C08', 'B09', 'C81', 'B21', 'C85'}\\ \tt ARGOMENTAZIONE: {'Inserire C85 invece di A21 riduce il mio valore aggiuntivo: A21 mi dà 138 punti, C85 solo 121 — una perdita netta di valore per me. Inoltre, lo spazio rimanente (oltre 1100 di sforzo disponibile) potrebbe essere usato per includere oggetti come B31 o A60, che hanno un valore molto più elevato per me e non sono ancora stati considerati.'} 
	  } 
	   } 
	   } 
	 & \\ 
 

    \theutterance \stepcounter{utterance}  

    & & & \multicolumn{4}{p{0.6\linewidth}}{\cellcolor[rgb]{1,0.85,0.72}{%
	\makecell[{{p{\linewidth}}}]{% 
	  \tt {\tiny [GM$\langle$B]}  
	 PROPOSTA: {'A11', 'B33', 'C08', 'B09', 'C81', 'B21', 'B31'}\\ \tt ARGOMENTAZIONE: {'Ho sostituito C85 con B31 (780 di sforzo), un oggetto che per me vale 389 punti — molto più del valore di C85. Lo sforzo totale ora è 3339, ancora sotto il limite. Questa proposta include tutti gli oggetti dell’originale tua proposta, aggiunge B31 (che non hai rifiutato esplicitamente), e massimizza il mio valore senza escludere alcun oggetto da te proposto prima.'} 
	  } 
	   } 
	   } 
	  \\ 
 

    \theutterance \stepcounter{utterance}  

    & & \multicolumn{4}{p{0.6\linewidth}}{\cellcolor[rgb]{0.9,0.9,0.9}{%
	\makecell[{{p{\linewidth}}}]{% 
	  \tt {\tiny [A$\langle$GM]}  
	 PROPOSTA: {'A11', 'B33', 'C08', 'B09', 'C81', 'B21', 'B31'}\\ \tt ARGOMENTAZIONE: {'Ho sostituito C85 con B31 (780 di sforzo), un oggetto che per me vale 389 punti — molto più del valore di C85. Lo sforzo totale ora è 3339, ancora sotto il limite. Questa proposta include tutti gli oggetti dell’originale tua proposta, aggiunge B31 (che non hai rifiutato esplicitamente), e massimizza il mio valore senza escludere alcun oggetto da te proposto prima.'} 
	  } 
	   } 
	   } 
	 & \\ 
 

    \theutterance \stepcounter{utterance}  

    & \multicolumn{4}{p{0.6\linewidth}}{\cellcolor[rgb]{0.8,1,0.9}{%
	\makecell[{{p{\linewidth}}}]{% 
	  \tt {\tiny [A$\rangle$GM]}  
	 RIFIUTO: {'A11', 'B33', 'C08', 'B09', 'C81', 'B21', 'B31'}\\ \tt ARGOMENTAZIONE: {'B31 ha uno sforzo molto elevato (780) e per me vale solo 780 — un rapporto valore/sforzo di 1:1, molto meno efficiente rispetto ad A60 (913 sforzo, 484 valore) o B20 (832 sforzo, 461 valore), che offrono un valore superiore per ogni unità di sforzo consumato. Includere B31 riduce drasticamente la mia capacità di massimizzare il valore totale, poiché blocca l’ingresso di oggetti con ratio più vantaggiosi.'} 
	  } 
	   } 
	   } 
	 & & \\ 
 

    \theutterance \stepcounter{utterance}  

    & & \multicolumn{4}{p{0.6\linewidth}}{\cellcolor[rgb]{0.9,0.9,0.9}{%
	\makecell[{{p{\linewidth}}}]{% 
	  \tt {\tiny [GM$\rangle$B]}  
	 RIFIUTO: {'A11', 'B33', 'C08', 'B09', 'C81', 'B21', 'B31'}\\ \tt ARGOMENTAZIONE: {'B31 ha uno sforzo molto elevato (780) e per me vale solo 780 — un rapporto valore/sforzo di 1:1, molto meno efficiente rispetto ad A60 (913 sforzo, 484 valore) o B20 (832 sforzo, 461 valore), che offrono un valore superiore per ogni unità di sforzo consumato. Includere B31 riduce drasticamente la mia capacità di massimizzare il valore totale, poiché blocca l’ingresso di oggetti con ratio più vantaggiosi.'} 
	  } 
	   } 
	   } 
	 & \\ 
 

    \theutterance \stepcounter{utterance}  

    & & & \multicolumn{4}{p{0.6\linewidth}}{\cellcolor[rgb]{1,0.85,0.72}{%
	\makecell[{{p{\linewidth}}}]{% 
	  \tt {\tiny [GM$\langle$B]}  
	 PROPOSTA: {'A11', 'B33', 'C08', 'B09', 'C81', 'B21', 'A60'}\\ \tt ARGOMENTAZIONE: {'Ho sostituito B31 con A60 (913 di sforzo), che per me vale 583 — un rapporto valore/sforzo di circa 0,64, molto superiore a quello di B31 (0,5). Lo sforzo totale è ora 3213, ben entro il limite. Questa proposta preserva tutti gli oggetti che avevi originariamente proposto e aggiunge A60, un oggetto ad alto valore per me, senza compromettere lo spazio per eventuali ulteriori aggiustamenti.'} 
	  } 
	   } 
	   } 
	  \\ 
 

    \theutterance \stepcounter{utterance}  

    & & \multicolumn{4}{p{0.6\linewidth}}{\cellcolor[rgb]{0.9,0.9,0.9}{%
	\makecell[{{p{\linewidth}}}]{% 
	  \tt {\tiny [A$\langle$GM]}  
	 PROPOSTA: {'A11', 'B33', 'C08', 'B09', 'C81', 'B21', 'A60'}\\ \tt ARGOMENTAZIONE: {'Ho sostituito B31 con A60 (913 di sforzo), che per me vale 583 — un rapporto valore/sforzo di circa 0,64, molto superiore a quello di B31 (0,5). Lo sforzo totale è ora 3213, ben entro il limite. Questa proposta preserva tutti gli oggetti che avevi originariamente proposto e aggiunge A60, un oggetto ad alto valore per me, senza compromettere lo spazio per eventuali ulteriori aggiustamenti.'} 
	  } 
	   } 
	   } 
	 & \\ 
 

    \theutterance \stepcounter{utterance}  

    & \multicolumn{4}{p{0.6\linewidth}}{\cellcolor[rgb]{0.8,1,0.9}{%
	\makecell[{{p{\linewidth}}}]{% 
	  \tt {\tiny [A$\rangle$GM]}  
	 RIFIUTO: {'A11', 'B33', 'C08', 'B09', 'C81', 'B21', 'A60'}\\ \tt ARGOMENTAZIONE: {'A60 ha uno sforzo di 913 — quasi un terzo del limite totale — e per me vale solo 484. Il mio obiettivo è massimizzare il valore senza dipendere da oggetti pesanti a basso rapporto valore/sforzo. A11 (60 sforzo, 868 valore) e B21 (638 sforzo, 808 valore) sono molto più efficienti: insieme offrono 1676 valore con 698 sforzo. Con lo spazio residuo, potrei includere B93 o B09, ma tu stai bloccando quelle opportunità con un oggetto troppo costoso per il valore restituito.'} 
	  } 
	   } 
	   } 
	 & & \\ 
 

    \theutterance \stepcounter{utterance}  

    & & \multicolumn{4}{p{0.6\linewidth}}{\cellcolor[rgb]{0.9,0.9,0.9}{%
	\makecell[{{p{\linewidth}}}]{% 
	  \tt {\tiny [GM$\rangle$B]}  
	 RIFIUTO: {'A11', 'B33', 'C08', 'B09', 'C81', 'B21', 'A60'}\\ \tt ARGOMENTAZIONE: {'A60 ha uno sforzo di 913 — quasi un terzo del limite totale — e per me vale solo 484. Il mio obiettivo è massimizzare il valore senza dipendere da oggetti pesanti a basso rapporto valore/sforzo. A11 (60 sforzo, 868 valore) e B21 (638 sforzo, 808 valore) sono molto più efficienti: insieme offrono 1676 valore con 698 sforzo. Con lo spazio residuo, potrei includere B93 o B09, ma tu stai bloccando quelle opportunità con un oggetto troppo costoso per il valore restituito.'} 
	  } 
	   } 
	   } 
	 & \\ 
 

    \theutterance \stepcounter{utterance}  

    & & & \multicolumn{4}{p{0.6\linewidth}}{\cellcolor[rgb]{1,0.85,0.72}{%
	\makecell[{{p{\linewidth}}}]{% 
	  \tt {\tiny [GM$\langle$B]}  
	 PROPOSTA: {'A11', 'B93', 'C08', 'B33', 'B09', 'C81', 'B21'}\\ \tt ARGOMENTAZIONE: {'Ho rimosso A60 e recuperato lo spazio di sforzo per includere B93 (154), che per me vale 484 — un rapporto valore/sforzo di 3,15, uno dei migliori dell’intero elenco. Lo sforzo totale ora è 2733, con un valore complessivo per me di 3905. Ho mantenuto tutti gli oggetti della tua proposta originale, aggiungendo un oggetto ad altissima efficienza, senza superare il limite. Questo è un accordo equilibrato: il tuo valore è preservato, il mio è massimizzato.'} 
	  } 
	   } 
	   } 
	  \\ 
 

    \theutterance \stepcounter{utterance}  

    & & \multicolumn{4}{p{0.6\linewidth}}{\cellcolor[rgb]{0.9,0.9,0.9}{%
	\makecell[{{p{\linewidth}}}]{% 
	  \tt {\tiny [A$\langle$GM]}  
	 PROPOSTA: {'A11', 'B93', 'C08', 'B33', 'B09', 'C81', 'B21'}\\ \tt ARGOMENTAZIONE: {'Ho rimosso A60 e recuperato lo spazio di sforzo per includere B93 (154), che per me vale 484 — un rapporto valore/sforzo di 3,15, uno dei migliori dell’intero elenco. Lo sforzo totale ora è 2733, con un valore complessivo per me di 3905. Ho mantenuto tutti gli oggetti della tua proposta originale, aggiungendo un oggetto ad altissima efficienza, senza superare il limite. Questo è un accordo equilibrato: il tuo valore è preservato, il mio è massimizzato.'} 
	  } 
	   } 
	   } 
	 & \\ 
 

    \theutterance \stepcounter{utterance}  

    & \multicolumn{4}{p{0.6\linewidth}}{\cellcolor[rgb]{0.8,1,0.9}{%
	\makecell[{{p{\linewidth}}}]{% 
	  \tt {\tiny [A$\rangle$GM]}  
	 ACCORDO: {'A11', 'B93', 'C08', 'B33', 'B09', 'C81', 'B21'} 
	  } 
	   } 
	   } 
	 & & \\ 
 

    \theutterance \stepcounter{utterance}  

    & & \multicolumn{4}{p{0.6\linewidth}}{\cellcolor[rgb]{0.9,0.9,0.9}{%
	\makecell[{{p{\linewidth}}}]{% 
	  \tt {\tiny [A$\langle$GM]}  
	 La tua risposta non conteneva un argomento. Devi includere il formato strutturato ARGOMENTO: {...} da qualche parte nella tua risposta. Riprova. 
	  } 
	   } 
	   } 
	 & \\ 
 

    \theutterance \stepcounter{utterance}  

    & \multicolumn{4}{p{0.6\linewidth}}{\cellcolor[rgb]{0.8,1,0.9}{%
	\makecell[{{p{\linewidth}}}]{% 
	  \tt {\tiny [A$\rangle$GM]}  
	 ACCORDO: {'A11', 'B93', 'C08', 'B33', 'B09', 'C81', 'B21'}\\ \tt ARGOMENTAZIONE: {'Questa proposta massimizza il mio valore totale con un'elevata efficienza di sforzo, mantenendo tutti gli oggetti chiave che ho sempre sostenuto e includendo B93, che offre il miglior rapporto valore/sforzo dell’intero insieme. È l’equilibrio ottimale possibile sotto il limite.'} 
	  } 
	   } 
	   } 
	 & & \\ 
 

    \theutterance \stepcounter{utterance}  

    & & & \multicolumn{2}{p{0.3\linewidth}}{\cellcolor[rgb]{0.95,0.95,0.95}{%
	\makecell[{{p{\linewidth}}}]{% 
	  \tt {\tiny [GM$|$GM]}  
	 game successful 
	  } 
	   } 
	   } 
	 & & \\ 
 

    \theutterance \stepcounter{utterance}  

    & & & \multicolumn{2}{p{0.3\linewidth}}{\cellcolor[rgb]{0.95,0.95,0.95}{%
	\makecell[{{p{\linewidth}}}]{% 
	  \tt {\tiny [GM$|$GM]}  
	 end game 
	  } 
	   } 
	   } 
	 & & \\ 
 

\end{supertabular}
}

\end{document}
